% ==============================================================================
% 3. LÀM SẠCH, TIỀN XỬ LÝ, ĐẢM BẢO CHẤT LƯỢNG, LƯU TRỮ
% ==============================================================================
\needspace{0.3\textheight}
\section{Quy trình làm sạch và tiền xử lý dữ liệu}
\label{sec:cleaning}

Nguyên tắc chung của nhóm là \emph{chỉ lọc dữ liệu bằng quy tắc, không lọc bằng mô hình}, và không chỉnh sửa dữ liệu gốc. Bảng~\ref{tab:cleaning} trình bày từng bước làm sạch, cách thực hiện và số lượng ảnh bị loại.

\begin{longtable}{>{\raggedright\arraybackslash}p{0.21\linewidth} >{\raggedright\arraybackslash}p{0.44\linewidth} >{\raggedright\arraybackslash}p{0.26\linewidth}}
\caption{Các bước làm sạch và tiền xử lý, theo khung báo cáo}
\label{tab:cleaning}\\
\toprule
\textbf{Nội dung} & \textbf{Cách thực hiện} & \textbf{Kết quả} \\
\midrule
\endfirsthead
\toprule
\textbf{Nội dung} & \textbf{Cách thực hiện} & \textbf{Kết quả} \\
\midrule
\endhead
Ảnh chụp ban đêm & Loại ảnh có độ sáng trung bình nửa trên của ảnh $360^\circ$ dưới 90 (thang 0--255). Không lọc theo giờ chụp vì giờ ghi trên Mapillary không đáng tin cậy (nhiều ảnh ghi 18--0 giờ vẫn sáng như ban ngày). & Bộ chính: loại 58/3.057. Kho ứng viên bổ sung: loại 38/5.074 \\
Ảnh không thuộc tuyến đường & Chỉ giữ ảnh cách đường $\le 20$ m và có hướng di chuyển dọc trục đường ($\le 30^\circ$); loại hẻm. & 81,3\% ảnh trong khu vực gắn được vào đường \\
Ảnh không đủ thông tin & Không lọc tự động. Người gán đặt nhãn \lab{I} kèm lý do; ảnh \lab{I} bị loại khỏi mô hình nhưng vẫn được giữ trong bảng nhãn. & 1.920/5.149 ảnh mang nhãn \lab{I} \\
Trùng lặp tuyệt đối & Mỗi ảnh $360^\circ$ chỉ được dùng một lần; so sánh mã băm MD5 của ảnh cắt. & Xem mục~\ref{sec:eda-quality} \\
Gần trùng & Lấy mẫu thưa: mỗi 25 m một vị trí, tối đa 2 ảnh mỗi vị trí và phải khác ngày chụp. Riêng kho ứng viên bổ sung có thêm bước so sánh mã băm sai phân (dHash)~\cite{krawetz2013dhash} giữa hai ngày chụp của cùng vị trí. & Kho ứng viên bổ sung: loại 1 cặp. Bộ chính: được đánh giá lần đầu trong báo cáo này (mục~\ref{sec:eda-quality}) \\
Dữ liệu thiếu & Metadata lấy từ API; ảnh thiếu tọa độ hoặc hướng di chuyển không gắn được vào đường nên đã bị loại ở bước gắn đường. Với nhãn, mọi lượt gán đều bắt buộc chọn nhãn chính. & Xem bảng~\ref{tab:inventory} \\
Chuẩn hóa, biến đổi & Chiếu ảnh $360^\circ$ về ảnh phối cảnh $640\times640$ điểm ảnh với cùng tham số cho mọi ảnh (mục~\ref{sec:pipeline}). Chưa chuẩn hóa giá trị điểm ảnh; bước này thuộc giai đoạn huấn luyện. & Mọi ảnh có cùng kích thước và góc nhìn \\
Chuyển đổi kiểu dữ liệu & Nhãn Label Studio (chuỗi tiếng Việt) $\rightarrow$ mã \lab{0}/\lab{1}/\lab{I}. Mã ảnh $360^\circ$ (số 64 bit) $\rightarrow$ mã ảnh 8 ký tự không để lộ nguồn cho người gán. Thời gian chụp được giữ dạng chuỗi ISO có múi giờ. & Áp dụng thống nhất cho mọi bảng nhãn của bộ chính, pilot và ảnh bổ sung \\
Kiểm tra nhất quán & Bước thu thập từ chối các đường đã thuộc tập khác; báo cáo này kiểm tra lại toàn bộ khóa, tập, đường và tọa độ. & Mục~\ref{sec:eda-quality} \\
\bottomrule
\end{longtable}

\textbf{Các việc chưa thực hiện:} rà soát nội dung nhạy cảm trong ảnh (ngoài phần làm mờ sẵn của Mapillary) và khóa danh sách chia tập.

% ------------------------------------------------------------------------------
\needspace{0.45\textheight}
\section{Đảm bảo chất lượng dữ liệu}
\label{sec:qa}

\subsection{Tiêu chí dữ liệu tốt cho bài toán này}

Bảng~\ref{tab:criteria} nêu sáu tiêu chí chất lượng và cách kiểm tra từng tiêu chí.

\begin{table}[H]
\centering
\caption{Tiêu chí chất lượng và cách kiểm tra}
\label{tab:criteria}
\begin{tabularx}{\textwidth}{>{\raggedright\arraybackslash}p{0.25\linewidth} L >{\raggedright\arraybackslash}p{0.21\linewidth}}
\toprule
\textbf{Tiêu chí} & \textbf{Cách kiểm tra} & \textbf{Loại kiểm tra} \\
\midrule
Nhãn đúng định nghĩa và nhất quán giữa người gán & Gán đôi val/test, gán chéo 15\% train, Cohen $\kappa$~\cite{cohen1960kappa} với ngưỡng 0,6 đặt ra trong đề cương; phân xử mọi ảnh bất đồng & Thủ công và tự động \\
Tập đánh giá theo phân phối tự nhiên & Val/test chỉ gồm ảnh bộ chính; ảnh được chọn theo điểm chỉ đưa vào train; ghi tầng lấy mẫu vào manifest & Theo thiết kế \\
Không rò rỉ giữa các tập & Chia theo đường; loại các vị trí cách val/test dưới 40 m khi bổ sung ảnh; kiểm tra khoảng cách và ảnh gần trùng giữa các tập (mục~\ref{sec:eda-quality}) & Tự động và kiểm tra trực quan \\
Truy được nguồn gốc & Mỗi ảnh có một dòng trong manifest: mã ảnh gốc, tọa độ, ngày chụp, tài khoản chụp, giấy phép, đường, tập, tầng lấy mẫu & Tự động \\
Ảnh quan sát được & Bộ lọc ảnh đêm; nhãn \lab{I} kèm lý do; thống kê độ sáng, độ nét (mục~\ref{sec:eda-image}) & Tự động và thủ công \\
Đủ mẫu của lớp hiếm & Đếm ảnh \lab{Bị chặn} ở từng tập, kèm khoảng tin cậy (mục~\ref{sec:eda-labels}) & Tự động \\
\bottomrule
\end{tabularx}
\end{table}

\subsection{Kiểm tra thủ công: đồng thuận và phân xử}

Mục này đánh giá mức nhất quán giữa người gán (bảng~\ref{tab:kappa}) và mô tả quy trình phân xử các ảnh bất đồng.

\begin{table}[H]
\centering
\caption{Đồng thuận giữa hai người gán trên bộ chính (Cohen $\kappa$)}
\label{tab:kappa}
\begin{tabularx}{\textwidth}{>{\raggedright\arraybackslash}p{0.22\linewidth} >{\raggedright\arraybackslash}p{0.2\linewidth} R R R R}
\toprule
\textbf{Nhóm ảnh} & \textbf{Cặp gán} & \textbf{Số ảnh} & \textbf{Khớp nhau} & \textbf{$\kappa$ 3 nhãn} & \textbf{$\kappa$ chỉ 0/1} \\
\midrule
Test & Học -- Loan & 674 & 79\% & 0,60 & 0,68 \\
Val & Minh -- Khải & 486 & 80\% & 0,64 & 0,55 \\
Train gán chéo & 10 cặp & 276 & 78\% & 0,61 & 0,56 \\
\bottomrule
\end{tabularx}
\note{$\kappa$ chỉ 0/1 được tính trên các ảnh mà cả hai người đều chọn 0 hoặc 1 (test 418 ảnh, val 211 ảnh).}
\end{table}

\begin{itemize}
  \item Cả ba nhóm trong bảng~\ref{tab:kappa} đều đạt ngưỡng 0,6 đặt ra trong đề cương; pilot (theo hướng dẫn v0) chỉ đạt Fleiss $\kappa$ 0,57. Khoảng tin cậy của các giá trị $\kappa$ \textbf{chưa được tính}.
  \item Trong 1.436 ảnh được gán đôi có 303 ảnh bất đồng: 193 ảnh \lab{0}/\lab{I}, 51 ảnh \lab{1}/\lab{I}, 59 ảnh \lab{0}/\lab{1}. Như vậy phần lớn bất đồng liên quan tới việc ảnh có đủ thông tin để đánh giá hay không.
  \item \textbf{Phân xử 303 ảnh} được thực hiện qua hai bước. (1) Trợ lý AI Claude (Anthropic) xem từng ảnh mà không thấy phiếu của người gán và đề xuất nhãn kèm lý do; sau đó xem lại riêng 21 ảnh có đề xuất khác cả hai phiếu, lần này có hiển thị phiếu, và sửa 10 ảnh. (2) Khải, người phân xử, duyệt cả 303 đề xuất, giữ nguyên 275 và sửa 28 (9,2\%). Nhãn cuối là nhãn do Khải quyết định. Cả 28 nhãn được sửa đều trùng với phiếu của một trong hai người gán ban đầu.
  \item \textbf{7 nhãn cuối khác cả hai phiếu.} Ở 7 ảnh, người phân xử giữ đề xuất của Claude dù đề xuất này khác cả hai phiếu: 3 ảnh test, 2 ảnh val và 2 ảnh train (ví dụ: hai phiếu là \lab{0} và \lab{I}, nhãn cuối là \lab{1}).
  \item \textbf{Phân xử chưa hoàn toàn độc lập.} Đề cương dự kiến người phân xử độc lập, nhưng Khải là một trong hai người gán ở 123/303 ảnh bất đồng, gồm \emph{cả 99 ảnh bất đồng của val} (val do Minh và Khải gán) và 24 ảnh train. Trong 123 ảnh này, nhãn cuối trùng với phiếu của Khải ở 37 ảnh (đề xuất của Claude trùng ở 32 ảnh); trong 14 nhãn Khải sửa, 9 nhãn được đổi về phiếu của chính Khải. Vì vậy, nhãn val cần được diễn giải có tính đến hạn chế này.
  \item \textbf{Pilot.} Cả 150 ảnh được xem lại theo hướng dẫn v1 theo cùng quy trình: Claude đề xuất, Phúc duyệt và sửa 3 ảnh. So với phiếu v0, 25 ảnh có nhãn cuối khác nhãn đa số; 5 ảnh khác hòa phiếu ở v0 nên không có nhãn đa số.
\end{itemize}

\textbf{Ảnh bổ sung.} Mỗi ảnh chỉ có một người gán nên không tính được $\kappa$. Ảnh được chia ngẫu nhiên cho người gán, nhưng tỉ lệ \lab{Bị chặn} của từng người chênh lệch rõ (Học 0,43, Loan 0,19, tính trên nhãn trước kiểm tra chéo; kiểm định $\chi^2$ trên bảng người gán $\times$ nhãn, $p = 4\cdot10^{-15}$). Giá trị $p$ này giả định các ảnh độc lập và chưa xét sự phụ thuộc theo đường, nên chỉ mang tính tham khảo; để xác nhận khác biệt giữa người gán, cần cho nhiều người cùng gán một mẫu ảnh độc lập. Một thành viên đã kiểm tra chéo độc lập 203 ảnh thuộc danh sách 264 ảnh bị nghi ngờ: 186 ảnh (92\%) trùng nhãn gốc và 17 nhãn được đổi. Do 203 ảnh này được chọn vì bị nghi ngờ, chúng \emph{không đại diện} cho 2.000 ảnh; vì vậy mức nhiễu nhãn của ảnh bổ sung \textbf{chưa đo được}.

Bảng~\ref{tab:label-source} tổng hợp nguồn gốc của nhãn cuối.

\begin{table}[H]
\centering
\caption{Nguồn của nhãn cuối}
\label{tab:label-source}
\begin{tabularx}{\textwidth}{>{\raggedright\arraybackslash}p{0.4\linewidth} r L}
\toprule
\textbf{Nhãn cuối đến từ} & \textbf{Ảnh} & \textbf{Ghi chú} \\
\midrule
Hai người gán trùng nhau & 1.133 & Val, test và phần train gán chéo của bộ chính \\
Một người gán (train bộ chính) & 1.563 & \\
Phân xử: giữ nhãn đề xuất & 275 & Khải giữ nhãn Claude đề xuất \\
Phân xử: sửa nhãn đề xuất & 28 & Khải sửa; trùng phiếu của một người gán \\
Xem lại pilot: giữ nhãn đề xuất & 147 & Phúc giữ nhãn Claude đề xuất \\
Xem lại pilot: sửa nhãn đề xuất & 3 & Phúc sửa \\
Một người gán (ảnh bổ sung) & 1.983 & 186 ảnh trong đó đã được kiểm tra chéo và trùng nhãn gốc \\
Kiểm tra chéo (ảnh bổ sung) & 17 & Nhãn được đổi theo người kiểm tra chéo \\
\bottomrule
\end{tabularx}
\end{table}

\subsection{Kiểm tra tự động}

\begin{itemize}
  \item Trong quy trình thu thập và gán nhãn: bước thu thập từ chối các đường đã thuộc tập khác và các đường trùng tên với đường val/test; bước cập nhật giao diện gán nhãn kiểm tra cả 5 dự án gán nhãn đều dùng đúng giao diện v1; bước đo đồng thuận kiểm tra đủ số lượt gán và phát hiện các lượt lưu trùng.
  \item Trong báo cáo này: kiểm tra khóa, tập, đường, tọa độ, ngoại lệ, ảnh trùng và gần trùng, khoảng cách giữa các tập; kết quả được trình bày ở mục~\ref{sec:eda-quality}.
\end{itemize}

% ------------------------------------------------------------------------------
\needspace{0.5\textheight}
\section{Lưu trữ và quản lý dữ liệu}
\label{sec:storage}

\begin{table}[H]
\centering
\caption{Tổ chức thư mục dữ liệu}
\label{tab:storage}
\begin{tabularx}{\textwidth}{>{\raggedright\arraybackslash}p{0.31\linewidth} L c}
\toprule
\textbf{Thư mục / file} & \textbf{Nội dung} & \textbf{Trong git} \\
\midrule
\f{configs/streets_v1.yaml} & 17 đường của bộ chính và tập của từng đường & Có \\
\f{configs/streets_train_extra.yaml} & 75 đường của đợt bổ sung & Có \\
\f{data/manifest/} & Mỗi ảnh một dòng: nguồn gốc, tọa độ, ngày chụp, đường, tập, tầng lấy mẫu; bảng phân công gán nhãn & Có \\
\f{data/labels/main_labels_v1.csv} & Nhãn cuối của bộ chính và pilot (3.149 dòng) & Có \\
\f{data/labels/train_extra_labels_v2.csv} & Nhãn cuối của ảnh bổ sung (2.000 dòng) & Có \\
\f{data/labels/*_report/} & Từng lượt gán, báo cáo đồng thuận, bảng phân xử & Có \\
\f{data/labels/raw/} & Bản xuất gốc từ Label Studio (có email tài khoản) & Không \\
\f{data/streets/} & Bảng xếp hạng đường, điểm của bộ chọn ảnh, kết quả phân tích thăm dò & CSV: có \\
\f{data/raw/} & OSM, metadata Mapillary, ảnh $360^\circ$, embedding & Không \\
\f{data/crops/} & Ảnh cắt $640\times640$, tách thư mục cho bộ chính, pilot và ảnh bổ sung & Không \\
\f{docs/data_report/} & Báo cáo này, hình và bảng số liệu EDA & Có \\
\bottomrule
\end{tabularx}
\end{table}

\begin{itemize}[nosep]
  \item \textbf{Khóa chính} là mã ảnh 8 ký tự. Bảng nhãn và manifest được ghép với nhau theo khóa này; manifest ghi đường dẫn tới ảnh cắt.
  \item \textbf{Ảnh không được đưa lên git} (bảng~\ref{tab:storage}) vì dung lượng và giấy phép; ảnh có thể được dựng lại từ manifest (mã ảnh gốc và góc quay) bằng script của nhóm. Ảnh dựng lại khác bản trên server do nén JPEG lại (lệch trung bình 1,35--1,5 mức xám, đo trên 5 ảnh).
  \item \textbf{Phiên bản.} Tên file bảng nhãn có số phiên bản; bản cũ (v1 của nhãn ảnh bổ sung, trước kiểm tra chéo) được giữ lại. Phiếu gốc của từng người vẫn lưu trên Label Studio; nhãn cuối chỉ lưu trong repo.
  \item \textbf{Thông tin xác thực} (token Mapillary, mật khẩu server) nằm trong file cấu hình riêng, không đưa lên git.
\end{itemize}
